% --- Archon physics formalization source begin ---
% archon:physics
% archon:covers PhyXMiniProblems/problem_phyx_mini_0770.lean
% archon:source-report reports/phyx_mini/problem_phyx_mini_0770.source.json

\chapter{Physics problem phyx\_mini\_0770}
\label{ch:PhyXMiniProblems_problem_phyx_mini_0770}

\paragraph{Problem source.}
\#\# Physical scenario

A pulley on a friction-less axle has the shape of a uniform solid disk of mass 2.50 kg and ra-dius 20.0 cm. A 1.50 kg stone is attached to a very light wire that is wrapped around the rim of the pul-ley as shown in figure, and the system is released from rest.

\#\# Question

How far must the stone fall so that the pulley has 4.50 J of kinetic energy?

\#\# Answer choices

- A: A: 0.712m

- B: B: 0.673m

- C: C: 0.364m

- D: D: 0.198m

\#\# Figure caption (auxiliary; use the image as primary evidence)

The image shows a simple mechanical system comprising a pulley and a stone. The pulley is illustrated as a large circle, labeled "2.50 kg pulley." A string or rope is attached to the pulley, extending downward to the stone. The stone is depicted at the bottom with a label "1.50 kg stone." The setup suggests a physics problem typically involving pulleys and weights, likely related to tension, forces, or motion. There are no other elements or background features present in the image.

\#\# Dataset metadata

Work and Energy; Physical Model Grounding Reasoning, Multi-Formula Reasoning

\paragraph{Recorded answer/context.}
B: B: 0.673m

\paragraph{Figure/image path.}
/root/proposal\_for\_physic/science-mango/phyx\_mini\_run/phyx\_data/test\_image/770.png

\paragraph{Formalization target.}
create a compiling Lean file with sorry bodies at `PhyXMiniProblems/problem\_phyx\_mini\_0770.lean`. The Lean declarations must preserve the physical quantities, dimensions or dimensional roles, figure labels, governing-law hypotheses, and final relation expressed by this problem.
Use Mathlib/Physlib names found through LeanExplore where available. If a physics API is missing, introduce faithful local abstractions rather than scalar placeholder aliases.

\begin{theorem}[Physics formalization target]
\label{thm:physics:phyx_mini_0770:target}
\lean{PhyXMiniProblems.ProblemPhyXMini0770.fallDistance_matches_answerB}
\uses{def:physics:phyx-mini-0770:phyxminiproblems-problemphyxmini0770-lengthinmeters, def:physics:phyx-mini-0770:phyxminiproblems-problemphyxmini0770-pulleystonesetup, def:physics:phyx-mini-0770:phyxminiproblems-problemphyxmini0770-matchesproblemandfigurereadouts, def:physics:phyx-mini-0770:phyxminiproblems-problemphyxmini0770-usesstandardearthgravity, def:physics:phyx-mini-0770:phyxminiproblems-problemphyxmini0770-hasphysicalparameters, def:physics:phyx-mini-0770:phyxminiproblems-problemphyxmini0770-satisfiesidealpulleystonelaws, def:physics:phyx-mini-0770:phyxminiproblems-problemphyxmini0770-fitsdisplayedfalldistance}
The assigned autoformalize agent should translate this physics problem into Lean declarations in the covered file, with theorem and lemma proof bodies written as `by sorry`.
\end{theorem}
\begin{proof}
This is an autoformalization task, not a proof task. Produce statements that can later be proved by the physics prover without weakening the physical model.
\end{proof}

% --- Archon declaration topology begin ---
\section{Declaration topology}
\begin{definition}[Mass Quantity]
  \label{def:physics:phyx-mini-0770:phyxminiproblems-problemphyxmini0770-massquantity}
  \lean{PhyXMiniProblems.ProblemPhyXMini0770.MassQuantity}
  A nonnegative physical mass
\end{definition}
\begin{proof}
  This is the declared model definition of Mass Quantity.
\end{proof}
\begin{definition}[Length Quantity]
  \label{def:physics:phyx-mini-0770:phyxminiproblems-problemphyxmini0770-lengthquantity}
  \lean{PhyXMiniProblems.ProblemPhyXMini0770.LengthQuantity}
  A nonnegative physical length
\end{definition}
\begin{proof}
  This is the declared model definition of Length Quantity.
\end{proof}
\begin{definition}[Speed Quantity]
  \label{def:physics:phyx-mini-0770:phyxminiproblems-problemphyxmini0770-speedquantity}
  \lean{PhyXMiniProblems.ProblemPhyXMini0770.SpeedQuantity}
  A nonnegative linear-speed magnitude
\end{definition}
\begin{proof}
  This is the declared model definition of Speed Quantity.
\end{proof}
\begin{definition}[Angular Speed Quantity]
  \label{def:physics:phyx-mini-0770:phyxminiproblems-problemphyxmini0770-angularspeedquantity}
  \lean{PhyXMiniProblems.ProblemPhyXMini0770.AngularSpeedQuantity}
  A nonnegative angular-speed magnitude; radians are dimensionless
\end{definition}
\begin{proof}
  This is the declared model definition of Angular Speed Quantity.
\end{proof}
\begin{definition}[Acceleration Quantity]
  \label{def:physics:phyx-mini-0770:phyxminiproblems-problemphyxmini0770-accelerationquantity}
  \lean{PhyXMiniProblems.ProblemPhyXMini0770.AccelerationQuantity}
  A nonnegative acceleration magnitude
\end{definition}
\begin{proof}
  This is the declared model definition of Acceleration Quantity.
\end{proof}
\begin{definition}[Moment Of Inertia Quantity]
  \label{def:physics:phyx-mini-0770:phyxminiproblems-problemphyxmini0770-momentofinertiaquantity}
  \lean{PhyXMiniProblems.ProblemPhyXMini0770.MomentOfInertiaQuantity}
  Axial moment of inertia, with dimension mass times length squared
\end{definition}
\begin{proof}
  This is the declared model definition of Moment Of Inertia Quantity.
\end{proof}
\begin{definition}[Energy Quantity]
  \label{def:physics:phyx-mini-0770:phyxminiproblems-problemphyxmini0770-energyquantity}
  \lean{PhyXMiniProblems.ProblemPhyXMini0770.EnergyQuantity}
  Physical energy, with dimension mass times length squared per time squared
\end{definition}
\begin{proof}
  This is the declared model definition of Energy Quantity.
\end{proof}
\begin{definition}[mass In Kilograms]
  \label{def:physics:phyx-mini-0770:phyxminiproblems-problemphyxmini0770-massinkilograms}
  \lean{PhyXMiniProblems.ProblemPhyXMini0770.massInKilograms}
  \uses{def:physics:phyx-mini-0770:phyxminiproblems-problemphyxmini0770-massquantity}
  Read a mass in SI kilograms
\end{definition}
\begin{proof}
  This is the declared model definition of mass In Kilograms.
\end{proof}
\begin{definition}[length In Meters]
  \label{def:physics:phyx-mini-0770:phyxminiproblems-problemphyxmini0770-lengthinmeters}
  \lean{PhyXMiniProblems.ProblemPhyXMini0770.lengthInMeters}
  \uses{def:physics:phyx-mini-0770:phyxminiproblems-problemphyxmini0770-lengthquantity}
  Read a length in SI metres
\end{definition}
\begin{proof}
  This is the declared model definition of length In Meters.
\end{proof}
\begin{definition}[speed In Meters Per Second]
  \label{def:physics:phyx-mini-0770:phyxminiproblems-problemphyxmini0770-speedinmeterspersecond}
  \lean{PhyXMiniProblems.ProblemPhyXMini0770.speedInMetersPerSecond}
  \uses{def:physics:phyx-mini-0770:phyxminiproblems-problemphyxmini0770-speedquantity}
  Read a speed magnitude in SI metres per second
\end{definition}
\begin{proof}
  This is the declared model definition of speed In Meters Per Second.
\end{proof}
\begin{definition}[angular Speed In Radians Per Second]
  \label{def:physics:phyx-mini-0770:phyxminiproblems-problemphyxmini0770-angularspeedinradianspersecond}
  \lean{PhyXMiniProblems.ProblemPhyXMini0770.angularSpeedInRadiansPerSecond}
  \uses{def:physics:phyx-mini-0770:phyxminiproblems-problemphyxmini0770-angularspeedquantity}
  Read an angular speed in radians per SI second
\end{definition}
\begin{proof}
  This is the declared model definition of angular Speed In Radians Per Second.
\end{proof}
\begin{definition}[acceleration In Meters Per Second Squared]
  \label{def:physics:phyx-mini-0770:phyxminiproblems-problemphyxmini0770-accelerationinmeterspersecondsquared}
  \lean{PhyXMiniProblems.ProblemPhyXMini0770.accelerationInMetersPerSecondSquared}
  \uses{def:physics:phyx-mini-0770:phyxminiproblems-problemphyxmini0770-accelerationquantity}
  Read an acceleration in SI metres per second squared
\end{definition}
\begin{proof}
  This is the declared model definition of acceleration In Meters Per Second Squared.
\end{proof}
\begin{definition}[moment Of Inertia In Kilogram Meters Squared]
  \label{def:physics:phyx-mini-0770:phyxminiproblems-problemphyxmini0770-momentofinertiainkilogrammeterssquared}
  \lean{PhyXMiniProblems.ProblemPhyXMini0770.momentOfInertiaInKilogramMetersSquared}
  \uses{def:physics:phyx-mini-0770:phyxminiproblems-problemphyxmini0770-momentofinertiaquantity}
  Read an axial moment of inertia in kg m²
\end{definition}
\begin{proof}
  This is the declared model definition of moment Of Inertia In Kilogram Meters Squared.
\end{proof}
\begin{definition}[energy In Joules]
  \label{def:physics:phyx-mini-0770:phyxminiproblems-problemphyxmini0770-energyinjoules}
  \lean{PhyXMiniProblems.ProblemPhyXMini0770.energyInJoules}
  \uses{def:physics:phyx-mini-0770:phyxminiproblems-problemphyxmini0770-energyquantity}
  Read an energy in joules, calibrated against Physlib's joule
\end{definition}
\begin{proof}
  This is the declared model definition of energy In Joules.
\end{proof}
\begin{definition}[Figure Object]
  \label{def:physics:phyx-mini-0770:phyxminiproblems-problemphyxmini0770-figureobject}
  \lean{PhyXMiniProblems.ProblemPhyXMini0770.FigureObject}
  The three physical objects visible in the supplied bitmap
\end{definition}
\begin{proof}
  This is the declared model definition of Figure Object.
\end{proof}
\begin{definition}[Figure Text Label]
  \label{def:physics:phyx-mini-0770:phyxminiproblems-problemphyxmini0770-figuretextlabel}
  \lean{PhyXMiniProblems.ProblemPhyXMini0770.FigureTextLabel}
  The two text labels printed in the supplied bitmap
\end{definition}
\begin{proof}
  This is the declared model definition of Figure Text Label.
\end{proof}
\begin{definition}[Pulley Mass Model]
  \label{def:physics:phyx-mini-0770:phyxminiproblems-problemphyxmini0770-pulleymassmodel}
  \lean{PhyXMiniProblems.ProblemPhyXMini0770.PulleyMassModel}
  How the pulley's mass is distributed
\end{definition}
\begin{proof}
  This is the declared model definition of Pulley Mass Model.
\end{proof}
\begin{definition}[Axle Resistance Model]
  \label{def:physics:phyx-mini-0770:phyxminiproblems-problemphyxmini0770-axleresistancemodel}
  \lean{PhyXMiniProblems.ProblemPhyXMini0770.AxleResistanceModel}
  Resistance at the pulley's fixed axle
\end{definition}
\begin{proof}
  This is the declared model definition of Axle Resistance Model.
\end{proof}
\begin{definition}[Wire Mass Model]
  \label{def:physics:phyx-mini-0770:phyxminiproblems-problemphyxmini0770-wiremassmodel}
  \lean{PhyXMiniProblems.ProblemPhyXMini0770.WireMassModel}
  The mass idealization of the connecting wire
\end{definition}
\begin{proof}
  This is the declared model definition of Wire Mass Model.
\end{proof}
\begin{definition}[Wire Pulley Contact]
  \label{def:physics:phyx-mini-0770:phyxminiproblems-problemphyxmini0770-wirepulleycontact}
  \lean{PhyXMiniProblems.ProblemPhyXMini0770.WirePulleyContact}
  Contact between the wrapped wire and the pulley rim
\end{definition}
\begin{proof}
  This is the declared model definition of Wire Pulley Contact.
\end{proof}
\begin{definition}[Supplied Pulley Stone Figure]
  \label{def:physics:phyx-mini-0770:phyxminiproblems-problemphyxmini0770-suppliedpulleystonefigure}
  \lean{PhyXMiniProblems.ProblemPhyXMini0770.SuppliedPulleyStoneFigure}
  \uses{def:physics:phyx-mini-0770:phyxminiproblems-problemphyxmini0770-figureobject, def:physics:phyx-mini-0770:phyxminiproblems-problemphyxmini0770-figuretextlabel}
  Qualitative and printed information transcribed from the primary image
\end{definition}
\begin{proof}
  This is the declared model definition of Supplied Pulley Stone Figure.
\end{proof}
\begin{definition}[Pulley Stone Setup]
  \label{def:physics:phyx-mini-0770:phyxminiproblems-problemphyxmini0770-pulleystonesetup}
  \lean{PhyXMiniProblems.ProblemPhyXMini0770.PulleyStoneSetup}
  \uses{def:physics:phyx-mini-0770:phyxminiproblems-problemphyxmini0770-massquantity, def:physics:phyx-mini-0770:phyxminiproblems-problemphyxmini0770-lengthquantity, def:physics:phyx-mini-0770:phyxminiproblems-problemphyxmini0770-speedquantity, def:physics:phyx-mini-0770:phyxminiproblems-problemphyxmini0770-angularspeedquantity, def:physics:phyx-mini-0770:phyxminiproblems-problemphyxmini0770-accelerationquantity, def:physics:phyx-mini-0770:phyxminiproblems-problemphyxmini0770-momentofinertiaquantity, def:physics:phyx-mini-0770:phyxminiproblems-problemphyxmini0770-energyquantity, def:physics:phyx-mini-0770:phyxminiproblems-problemphyxmini0770-pulleymassmodel, def:physics:phyx-mini-0770:phyxminiproblems-problemphyxmini0770-axleresistancemodel, def:physics:phyx-mini-0770:phyxminiproblems-problemphyxmini0770-wiremassmodel, def:physics:phyx-mini-0770:phyxminiproblems-problemphyxmini0770-wirepulleycontact, def:physics:phyx-mini-0770:phyxminiproblems-problemphyxmini0770-suppliedpulleystonefigure}
  The independent physical quantities describing release and the later state at which the pulley's energy is inspected. In particular, fallDistance is a response field constrained by the governing laws below; it is not defined from any answer choice
\end{definition}
\begin{proof}
  This is the declared model definition of Pulley Stone Setup.
\end{proof}
\begin{definition}[Matches Problem And Figure Readouts]
  \label{def:physics:phyx-mini-0770:phyxminiproblems-problemphyxmini0770-matchesproblemandfigurereadouts}
  \lean{PhyXMiniProblems.ProblemPhyXMini0770.MatchesProblemAndFigureReadouts}
  \uses{def:physics:phyx-mini-0770:phyxminiproblems-problemphyxmini0770-massinkilograms, def:physics:phyx-mini-0770:phyxminiproblems-problemphyxmini0770-lengthinmeters, def:physics:phyx-mini-0770:phyxminiproblems-problemphyxmini0770-speedinmeterspersecond, def:physics:phyx-mini-0770:phyxminiproblems-problemphyxmini0770-angularspeedinradianspersecond, def:physics:phyx-mini-0770:phyxminiproblems-problemphyxmini0770-energyinjoules, def:physics:phyx-mini-0770:phyxminiproblems-problemphyxmini0770-pulleystonesetup}
  Numerical problem data and primary-image readouts. The radius is stated in the text but is not printed in the bitmap. Release from rest is stated in the text and is recorded for both translational and rotational motion
\end{definition}
\begin{proof}
  This is the declared model definition of Matches Problem And Figure Readouts.
\end{proof}
\begin{definition}[Uses Standard Earth Gravity]
  \label{def:physics:phyx-mini-0770:phyxminiproblems-problemphyxmini0770-usesstandardearthgravity}
  \lean{PhyXMiniProblems.ProblemPhyXMini0770.UsesStandardEarthGravity}
  \uses{def:physics:phyx-mini-0770:phyxminiproblems-problemphyxmini0770-accelerationinmeterspersecondsquared, def:physics:phyx-mini-0770:phyxminiproblems-problemphyxmini0770-pulleystonesetup}
  The conventional near-Earth gravitational acceleration used in the multiple-choice numerical evaluation
\end{definition}
\begin{proof}
  This is the declared model definition of Uses Standard Earth Gravity.
\end{proof}
\begin{definition}[Has Physical Parameters]
  \label{def:physics:phyx-mini-0770:phyxminiproblems-problemphyxmini0770-hasphysicalparameters}
  \lean{PhyXMiniProblems.ProblemPhyXMini0770.HasPhysicalParameters}
  \uses{def:physics:phyx-mini-0770:phyxminiproblems-problemphyxmini0770-massinkilograms, def:physics:phyx-mini-0770:phyxminiproblems-problemphyxmini0770-lengthinmeters, def:physics:phyx-mini-0770:phyxminiproblems-problemphyxmini0770-accelerationinmeterspersecondsquared, def:physics:phyx-mini-0770:phyxminiproblems-problemphyxmini0770-energyinjoules, def:physics:phyx-mini-0770:phyxminiproblems-problemphyxmini0770-pulleystonesetup}
  Positivity and nondegeneracy conditions for the physical experiment
\end{definition}
\begin{proof}
  This is the declared model definition of Has Physical Parameters.
\end{proof}
\begin{definition}[Satisfies Ideal Pulley Stone Laws]
  \label{def:physics:phyx-mini-0770:phyxminiproblems-problemphyxmini0770-satisfiesidealpulleystonelaws}
  \lean{PhyXMiniProblems.ProblemPhyXMini0770.SatisfiesIdealPulleyStoneLaws}
  \uses{def:physics:phyx-mini-0770:phyxminiproblems-problemphyxmini0770-massinkilograms, def:physics:phyx-mini-0770:phyxminiproblems-problemphyxmini0770-lengthinmeters, def:physics:phyx-mini-0770:phyxminiproblems-problemphyxmini0770-speedinmeterspersecond, def:physics:phyx-mini-0770:phyxminiproblems-problemphyxmini0770-angularspeedinradianspersecond, def:physics:phyx-mini-0770:phyxminiproblems-problemphyxmini0770-accelerationinmeterspersecondsquared, def:physics:phyx-mini-0770:phyxminiproblems-problemphyxmini0770-momentofinertiainkilogrammeterssquared, def:physics:phyx-mini-0770:phyxminiproblems-problemphyxmini0770-energyinjoules, def:physics:phyx-mini-0770:phyxminiproblems-problemphyxmini0770-pulleystonesetup}
  Governing laws For an ideal fixed-axis uniform disk and a taut non-slipping wire: I = M R² / 2; the rim speed equals the stone speed, v = R ω; the pulley's rotational kinetic energy is I ω² / 2; with a massless wire and frictionless axle, the loss m g h of gravitational potential energy equals the total gain in translational and rotational kinetic energy. These laws relate independent response fields. They contain neither the exact solved distance 33/49 m nor the displayed answer 0.673 m
\end{definition}
\begin{proof}
  This is the declared model definition of Satisfies Ideal Pulley Stone Laws.
\end{proof}
\begin{definition}[Answer Choice]
  \label{def:physics:phyx-mini-0770:phyxminiproblems-problemphyxmini0770-answerchoice}
  \lean{PhyXMiniProblems.ProblemPhyXMini0770.AnswerChoice}
  The four fall distances displayed in the problem
\end{definition}
\begin{proof}
  This is the declared model definition of Answer Choice.
\end{proof}
\begin{definition}[displayed Fall Distance Meters]
  \label{def:physics:phyx-mini-0770:phyxminiproblems-problemphyxmini0770-displayedfalldistancemeters}
  \lean{PhyXMiniProblems.ProblemPhyXMini0770.displayedFallDistanceMeters}
  \uses{def:physics:phyx-mini-0770:phyxminiproblems-problemphyxmini0770-answerchoice}
  Numerical metre readout printed beside an answer choice
\end{definition}
\begin{proof}
  This is the declared model definition of displayed Fall Distance Meters.
\end{proof}
\begin{definition}[Fits Displayed Fall Distance]
  \label{def:physics:phyx-mini-0770:phyxminiproblems-problemphyxmini0770-fitsdisplayedfalldistance}
  \lean{PhyXMiniProblems.ProblemPhyXMini0770.FitsDisplayedFallDistance}
  \uses{def:physics:phyx-mini-0770:phyxminiproblems-problemphyxmini0770-lengthinmeters, def:physics:phyx-mini-0770:phyxminiproblems-problemphyxmini0770-pulleystonesetup, def:physics:phyx-mini-0770:phyxminiproblems-problemphyxmini0770-answerchoice, def:physics:phyx-mini-0770:phyxminiproblems-problemphyxmini0770-displayedfalldistancemeters}
  A displayed three-decimal distance represents the modeled distance when the discrepancy is strictly less than half of the last displayed metre digit
\end{definition}
\begin{proof}
  This is the declared model definition of Fits Displayed Fall Distance.
\end{proof}
\begin{lemma}[fall Distance exact]
  \label{lem:physics:phyx-mini-0770:phyxminiproblems-problemphyxmini0770-falldistance-exact}
  \lean{PhyXMiniProblems.ProblemPhyXMini0770.fallDistance_exact}
  \uses{def:physics:phyx-mini-0770:phyxminiproblems-problemphyxmini0770-lengthinmeters, def:physics:phyx-mini-0770:phyxminiproblems-problemphyxmini0770-pulleystonesetup, def:physics:phyx-mini-0770:phyxminiproblems-problemphyxmini0770-matchesproblemandfigurereadouts, def:physics:phyx-mini-0770:phyxminiproblems-problemphyxmini0770-usesstandardearthgravity, def:physics:phyx-mini-0770:phyxminiproblems-problemphyxmini0770-hasphysicalparameters, def:physics:phyx-mini-0770:phyxminiproblems-problemphyxmini0770-satisfiesidealpulleystonelaws}
  The ideal-energy model determines the exact fall distance h = 33/49 m
\end{lemma}
\begin{proof}
  The conclusion follows from the governing relations and figure readouts named in the statement of fall Distance exact.
\end{proof}
% --- Archon declaration topology end ---
% --- Archon physics formalization source end ---
